\section{Conclusion}
\label{sec:conclusion}
% [2 pts] Wrap up all the context in your report and give a short summary.

We implemented SDS, VSD, and SDI for latent image optimization using Stable Diffusion 2.1-base.
SDI achieved the highest average CLIP score of 0.3086, followed by SDS at 0.3042 and VSD at 0.3011.
SDS produced recognizable but smoothed subjects, VSD recovered more detail with irregular textures,
and SDI generally produced clearer structures and cleaner surfaces.
VSD uses a learned noise predictor adapted to the current samples, while SDI obtains image-dependent
perturbations through inversion, providing alternatives to the sampled-noise reference in SDS.
These changes can improve visual detail, although VSD did not exceed SDS's average CLIP score.

The ablations showed that stronger SDS guidance helps only up to a point, additional inversion steps
do not guarantee higher scores, and repeatedly refreshing the SDI target is preferable to fitting
only its initial estimate.
Directly optimizing noise-prediction MSE through the U-Net produced noisy, unrecognizable images,
further illustrating the importance of the chosen update direction.
Overall, the results highlight the influence of noise construction and target updates on generation
quality, and the need to assess both text alignment and visual appearance.
